% Appendix D: exactness contracts and the stock engine.
% Definitions, Proposition gdn: proved. Contract witnesses: register entry contracts.
% Divergence mechanisms, noise floor, tap controls: register entries state-mech,
% state-noise, state-tap, state-pinned. Cache-level checks (plain vs MTP, the two
% history pairs, prefix cache on vs off, one prompt's same-server repeats): state-cache.
% Cache-level check of the concurrency pair and of the other four tap-changed prompts,
% remaining pairs, remaining H5 tests: pending.
% Figure fig:divergence: counts typed from evidence/state_safety/mechanism_*.json
% (summary.first_difference_module); switch to first_difference_by_module.csv when it
% is committed.
\section{Exactness contracts and the stock engine}\label{app:exactness}\label{app:engine}
Section~\ref{sec:contract} names the stock head kernel at the same batch shape as the reference for the certified head. This appendix places that choice among the contracts a decoder can promise, lists what each consumer of the head needs, and gives the evidence behind the same-shape rule: where stock configurations that ought to agree already disagree, and how the engine keeps its recurrent state under speculation.

\subsection{Five contracts and the witnesses that separate them}\label{app:contracts}\label{sec:contracts-p1}
An exactness claim needs a reference and a sense of equality. The reference is an executable specification, not an idealized network: the weights, the kernels, the decoding policy, the random-number protocol and the execution shape, together with whether it is stock plain decoding, stock speculative verification or a recorded execution. Mathematical equivalence over the reals does not imply that two GPU reductions choose the same token, because a different batch shape or reduction order can move logits that are close~\cite{floatguide,torchnumeric}; a certificate must therefore enclose the value the reference computes, casts included. Table~\ref{tab:contracts} lists the contracts a change can claim against such a reference and the results of this paper that rest on each. They are ordered roughly from strongest to weakest, but none implies the next in general.

\begin{table}[htbp]
\small\centering
\begin{tabularx}{\linewidth}{@{}P{2.6cm}P{4.4cm}YP{2.9cm}@{}}\toprule
Contract & What must be equal & Results that claim it & Status\\\midrule
Identical internal state & Named logical tensors (recurrent state, attention cache) at named boundaries, bitwise; not layouts or scratch buffers & Replay of recorded GDN updates (Proposition~\ref{prop:gdn}); rounding-preserving replay (P4, Appendix~\ref{app:minimal-state}); SGLang's speculative state path (H5, Appendix~\ref{app:state}) & Proved; P4 and H5 tests \pend{moonshot-p4,state-targeted}\\
Identical greedy outputs & Decisions, ties and stopping, with a valid state from which to continue & Certified head under \Rstock{} (Theorem~\ref{thm:stock}, Proposition~\ref{prop:stockinterval}) & Proved under an error model; measured on replayed rows, through the kernel on 60,000 of them~\ev{kernel-replay}\\
Identical seeded samples & The reference's mapping from seed to output; with an inverse-CDF reference, certify with CDF bounds against the same uniform and fall back without redrawing & Certified race with the stock seeded field; fixed-noise verification against the seeded sampler on the verify pass's logits (Section~\ref{sec:decisions}) & Proved; on 60,000 replayed rows at $T=0.7$ and 1.0, without top-$k$ or top-$p$, the certified sampler with its fallback equals the stock seeded sampler on every row~\ev{kernel-seeded}; in the engine \pend{kernel-seeded-engine}\\
Exact sampling distribution & A named law: the ideal model's, the categorical law of the reference's scores, or the finite random-number implementation's; speculative correction targets the verifier's own $p$ & Races (Theorem~\ref{thm:race}); fixed-noise verification in law; residual races (Appendix~\ref{app:residual}) & Proved\\
Empirical quality equivalence & Task scores and distances within declared margins, with confidence intervals & The lossy stack under its quality budget (Appendix~\ref{app:stacks}) & Budget declared; measurement \pend{moonshot-frontiers}\\\bottomrule
\end{tabularx}
\caption{Exactness contracts and the results that claim them. ``Measured on replayed rows'' means captured engine states replayed offline, not a served run.}
\label{tab:contracts}
\end{table}

Small exact witnesses separate the contracts~\ev{contracts}. Reassociating an FP32 sum flips a greedy winner: $1+2^{-24}+2^{-24}$ summed left to right returns $1$ and ties a competing logit of $1$, so the first-index rule picks the competitor, while the regrouped sum returns $1+2^{-23}$ and wins, as it does in real arithmetic. Distinct FP32 values, $1$ and $1+2^{-10}$, fall in one BF16 rounding cell, so a BF16-output head ties them and the index rule decides. Logits $(1,0)$ and $(2,0)$ have the same greedy decision but different laws: $(\tfrac23,\tfrac13)$ against $(\tfrac45,\tfrac15)$ in the base-2 mass model, and $0.731$ against $0.881$ on the first token at temperature one. Two inverse-CDF samplers that list the tokens of $p=(\tfrac12,\tfrac12)$ in different orders have the same law but map the uniform $\tfrac14$ to different tokens. Evaluating the gated-delta recurrence $S_t=S_{t-1}A_t+B_t$ in the split form $S_t=S_0P_t+W_t$, which is equal in exact arithmetic, changed 1,481 of 1,600 FP32 outputs over 50 steps at width 32, and 2 of them after rounding to BF16. A state difference the current query cannot see becomes visible after one recurrent step (Appendix~\ref{app:minimal-state}), so equal outputs now do not imply equal state. Two greedy decoders that each answer two of four problems correctly have equal accuracy, yet their answers can differ on every problem, so empirical quality equivalence implies none of the stronger contracts. Equal token sequences, finally, say nothing about latency or about whether the cache that produced them is owned correctly. A last witness, in the exact floating-point reference, separates \Rbf{} from \Rstock{}: with $w_0=(16,0,0)$, $w_1=(16,\tfrac1{16},2^{-30})$ and $h=(1,1,1)$, the real logits are $16$ and $16.0625+2^{-30}$, so \Rbf{} rounds them to $16$ and $16.125$ and picks token~1, while a sequential FP32 accumulation loses the $2^{-30}$, rounds $16.0625$ to even and returns $16$ for both rows, so the stock argmax picks token~0; the gap condition refuses to certify this row~\ev{precision}. The floating-point witnesses use IEEE binary32 with round to nearest and every reduction written as an explicit sequential loop; a kernel that contracts, reorders or accumulates differently can round differently. They show that the contracts are distinct in general and say nothing about how often the differences occur in the engine.

\subsection{The commit rule in general}\label{app:commit}\label{sec:commit}
Every certificate in this paper instantiates one rule. Partial computation yields information $I$ about the final logits: an approximate pass with error bounds, a transported summary, or an approximate hidden state with a radius. From $I$ we keep a \emph{sound set} $Z(I)\subseteq\R^V$ that contains the reference implementation's logits. A decision rule $D(z,R)$ maps logits and the randomness fixed by the contract (a noise field, a uniform draw) to a token or an accept bit.
\begin{definition}[Commit rule]\label{def:commit}
Commit the decision $d$ when $\{D(z,R):z\in Z(I)\}=\{d\}$. Otherwise acquire more information, which must shrink or replace $Z(I)$ by another sound set, and test again.
\end{definition}
If the reference logits $z^\star$ lie in $Z(I)$, a committed $d$ equals $D(z^\star,R)$; that is the whole soundness argument, and it holds whatever produced $Z(I)$. For a box $Z(I)=\prod_j[l_j,u_j]$ the greedy rule commits $c$ when $l_c>\max_{j\neq c}u_j$, and a Gumbel race with a fixed field $G$ commits $c$ when $l_c/T+G_c>\max_{j\ne c}(u_j/T+G_j)$, which needs no partition function as long as the field stays fixed through every refinement~\cite{astar}. When the information is about the hidden state rather than the logits, a box is loose and the directional certificate of Proposition~\ref{prop:directional} is the right one. The rule names the reference; it is the gap condition of Theorem~\ref{thm:stock} that turns a sound set for the real logits into one for the stock kernel's decision.

\subsection{What each consumer of the head needs}\label{app:consumers}\label{sec:consumers}
Table~\ref{tab:consumers} lists the information each consumer of the head requires. A fast path is admitted by these contracts, not by a request being labelled ``decode''. Grammar masks, repetition penalties and temperature transforms are part of the scores; they cannot be applied after candidates have been discarded.

\begin{table}[htbp]
\small\centering
\begin{tabularx}{\linewidth}{@{}P{3.4cm}YP{5.2cm}@{}}\toprule
Consumer & Required information & Certificate in this paper\\\midrule
Greedy token & Winning index under the reference kernel and tie rule & Ladder plus stock contract (Theorems~\ref{thm:ladder} and~\ref{thm:stock})\\
Exact top-$k$ & Winning indices in order, and their scores; no partition function & Ladder on membership and on the order within the set (Appendix~\ref{app:truncated})\\
Seeded or fixed-noise sample, full vocabulary & Argmax of perturbed scores under the token-keyed field & Certified race (Theorem~\ref{thm:race})\\
Seeded top-$k$ or top-$p$ sample (stock) & Each candidate's rank in the sorted probabilities, which keys its noise, and the sort's order on ties; for top-$p$, the mass above each token & Membership and complete order within the set, then the race (Appendix~\ref{app:truncated}; analysis)\\
Acceptance bit & $\exp(z_X)$ against $Uq_XZ$; an interval can suffice & Guards (Appendix~\ref{app:guards}; analysis). The engine path uses fixed-noise verification instead\\
Residual sample & The $(p-q)_+$ law, including omitted mass & Residual races (Appendix~\ref{app:residual}; analysis)\\
Bonus sample & A sample from the target law under the random-number contract & Certified race\\
Top-$p$ sample & Mass above each token relative to $Z$ & Partition bounds, or the ordinary path\\
Arbitrary log-probabilities & Requested logits and the full partition function & None; the ordinary path\\\bottomrule
\end{tabularx}
\caption{Consumers of the head and what each needs. The last rows are where a narrow fast path usually loses.}
\label{tab:consumers}
\end{table}

\subsection{Where equivalent stock configurations disagree}\label{app:divergence}\label{sec:divergence}
Plain decoding at different batch sizes, plain decoding against MTP speculation, and speculation at different batch sizes can emit different greedy tokens. Both greedy paths reduce to one rule, \code{torch.argmax} (first maximal index) over the BF16 head output widened exactly to FP32, so two runs that feed identical logits to it cannot choose different tokens. Every divergence therefore means that the computed logits differed, and the question is where. ``A tie'' is an observation about where the divergence became visible, not an explanation of why it happened. Batch-invariant kernels make greedy outputs reproducible across batch compositions~\cite{he2025nondeterminism}, and floating-point nondeterminism in inference has been characterized and partly mitigated~\cite{yuan2025nondeterminism}; what follows locates the kernels responsible in this engine. In brief: neither the tie rule nor the head GEMM caused a divergence~\ev{state-mech}; plain decoding and MTP speculation part inside layer 0's recurrence, whose decode and verify kernels produce different bits from identical inputs and state~\ev{state-cache}; plain decoding at two concurrencies parts at batch-dependent kernels~\ev{state-mech}; with identical pools, MTP diverges from plain decoding at 3.5 to 4.0 per 1,000 compared tokens, always at an exact tie or within two BF16 spacings~\ev{state-pinned}; and with the prefix cache on, a request's output can depend on which earlier request computed its shared prefix and on what the cache holds when it is prefilled~\ev{targeted,state-cache}.

\paragraph{Differential tests} A fixed set of 320 prompts from GSM8K, HumanEval, MT-Bench, AlpacaEval and CNN/DailyMail at pinned revisions, tokenized once so that every configuration receives identical token IDs (105 of them in the model's thinking mode), is decoded greedily for up to 256 tokens with the top five log-probabilities at every position. The configurations differ only in settings that should leave greedy output unchanged: plain decoding against MTP speculation (one, three and five draft steps with top-$k$ 1, and a small tree with top-$k$ 2), the prefix cache and the overlap scheduler on or off, \code{-{}-enable-deterministic-inference} and \code{-{}-enable-fp32-lm-head}, each with one request in flight and with a client concurrency of 32. In the first matrix every server was capped at 16 running requests, so ``concurrency 32'' meant decode batches of up to 16 with the rest queued; the pinned matrix below caps every server at 8 with identical pools. Two runs of a prompt are compared position by position; the first position where the tokens differ is a divergence, the positions before it are the exposure, and the rate per 1,000 tokens is divergences over exposure. Every run uses the engine at the pin with CUDA graphs, the overlap scheduler and the FlashInfer attention backend~\ev{state-noise}.

\paragraph{Locating the first difference} A diagnostic tap (\path{engine/sglang/patches/state/0001-state-tap.patch}) hashes the exact bits of every module output of the target, one row per token, and two points inside the GDN backend (after the causal convolution and after the recurrent kernel), into static device buffers that the CUDA graphs capture and replay; for selected requests it also saves the head input and the logits fed to the argmax. For each divergence both configurations are replayed on the prompt and the committed rows (for MTP, only verify rows on the accepted path) are walked in order to the first position and, in execution order, the first module whose output bits differ. Before that point every hashed module output agrees. The caches (attention KV, the GDN convolution window and the recurrent state) are not module outputs, so a difference written into a cache without a differing output would first show in the module that reads it. A later version of the tap closes that gap: it also hashes the caches each forward reads, the keys and values of every attention layer at every cached position and the convolution window and recurrent state of every GDN layer~\ev{state-cache}. The comparison skips the caches of a prefill with no cached prefix, which reads no earlier state: its recurrent slot can still hold an earlier request's state when the tap hashes it, because the engine clears the slot inside the forward. This version has been run at batch 1 for plain decoding against MTP, for the history pairs, for the prefix cache on against off and for one prompt's same-server repeats (below); for the concurrency pair the check is \pend{state-mech}. The tool was checked on itself: a one-ulp change injected into layer 9's MLP down-projection output is named at that module, at the first prompt token, in 4 of 4 prompts, and one injected into layer 13's GDN recurrence output is named at that operation in 4 of 4 when the search starts at the first decode position; searched from the prompt, the second is attributed to the enclosing GDN attention module, because in prefill the attention backends run eagerly between captured graph segments~\ev{state-tap}. With the tap version behind Figure~\ref{fig:divergence}, tapped runs match untapped ones in tokens and log-probabilities for 162 of 167 prompts~\ev{state-mech}, and with the cache-hashing version for 35 of 40, the other five being the same five prompts~\ev{state-cache}. The tapped and untapped servers allocated different pools, 159,322 (159,324 for the cache-hashing version) against 97,672 attention-cache tokens and 199 against 122 recurrent-state slots, so the five exceptions are not attributed to the tap itself~\ev{state-pools}; they are discussed below, with the history dependence through the prefix cache. Outside the pinned matrix below, every comparison in this section across two servers had different pools, while the concurrency comparison, the same-server repeats and the deterministic-mode checks compare passes on one server.

\paragraph{Findings that do not depend on an error model} Two pairs have been analysed~\ev{state-mech}. First, neither the tie rule nor the head GEMM caused a single divergence: in every case the head input itself already differed. Second, the first differing module output is specific to each pair (Figure~\ref{fig:divergence}). For plain decoding against MTP speculation with three draft steps, both at batch 1, it is layer 0's GDN recurrence output at the first speculative cycle, in all 167 prompts analysed, while that layer's causal-convolution output, the recurrence's immediate input, is bitwise equal. At the pin, decode and target verify call different Triton kernels (\code{fused\_recurrent\_gated\_delta\_rule\_packed\_decode} and \code{fused\_sigmoid\_gating\_delta\_rule\_update}), which fuse the gating and the state update differently. Rerun with the cache-hashing tap on 40 of these prompts, every cache entering every forward is identical up to that output, and the first cache to differ is layer 0's recurrent state, which the differing kernel writes~\ev{state-cache}; the two kernels therefore produce different bits from identical inputs and identical state. The two tapped servers of the 167-prompt analysis had different pools (159,322 against 166,522 attention-cache tokens and 199 against 97 recurrent-state slots; 159,324 against 166,481 tokens and 199 against 97 slots for the cache-hashing rerun), and MTP's draft layer attends over its own cached prefix, so a different copy of a shared prefix could change the proposals, the accepted lengths and the later verify forwards without any target attention output differing first. The location of the first difference bounds the pools only for that difference, which is at output index 1, the first row of the first verify forward, whose layer-0 recurrence output depends on its token and on the state left by the prefill, not on the proposals or on cached attention values (reasoning from the row-wise structure of the kernels); whether the pools contributed to the divergences that follow is \pend{state-mech}. For plain decoding at concurrency 1 against 32 (batches of up to 16; 40 tapped prompts), it is a GDN gated RMSNorm in decode in 29 prompts, a full-attention layer's output in decode in 6, and the MLP down-projection of a prefill batched with other requests in 5; never the GDN recurrence and never a decode GEMM. These two tapped servers also had different pools (159,322 against 89,652 attention-cache tokens and 199 against 111 recurrent-state slots). A different copy of a shared prefix's attention values would first show at a full-attention output, so for the 6 prompts that first differ there the pools and the prefix cache's history are possible causes besides the attention kernel's dependence on the batch; for the other 34, every attention output before the first difference agrees. The gated norm chooses how many rows each program reduces from the batch's row count (one row per program for a batch of 1, two for a batch of 16), which changes the reduction tile and so the FP32 sum order of a row; the same function returns a constant in batch-invariant mode. FlashInfer's decode plan and cuBLAS's GEMM choice also depend on the batch. For this pair these are first differing module \emph{outputs} only; that both runs enter the differing kernel with equal caches is \pend{state-mech}.

\begin{figure}[htbp]
\centering
% Figure: where equivalent stock configurations first produce different module outputs, by
% decoder layer. Drawn from evidence/state_safety/first_difference_by_module.csv (one row per
% pair, first differing module, layer, kind and count). Full-attention layers are every fourth
% layer (3, 7, ..., 31); the rest are Gated DeltaNet layers. The module kinds have no paper-wide
% colour, so every mark is ink and the kinds differ by marker shape; each mark's count is printed.
\datafigure{../evidence/state_safety/first_difference_by_module.csv}{%
\vpreadcsv{\divtable}{../evidence/state_safety/first_difference_by_module.csv}%
\vpsum{\divtable}{pair}{plain_c1_vs_mtp_s3_c1}{count}{\nmtp}%
\vpsum{\divtable}{pair}{plain_c1_vs_c32}{count}{\nconc}%
\begin{tikzpicture}[font=\footnotesize]
\begin{axis}[x=0.33cm,y=1cm,xmin=-0.6,xmax=31.6,ymin=-0.55,ymax=2.05,clip=false,
  axis x line=bottom,axis y line=none,axis line style={ink!60,-},
  xtick={0,4,8,12,16,20,24,28,31},tick align=outside,major tick length=0.1cm,typeset ticklabels with strut,
  tick style={ink!60},
  xlabel={decoder layer (shaded: full attention; others: Gated DeltaNet)},
  visualization depends on={value \thisrow{count}\as\cnt}]
% full-attention layer columns, above the axis line (no macro parameter: the figure is an
% argument of \datafigure)
\fill[wash] (axis cs:2.55,-0.5) rectangle (axis cs:3.45,2.05);\fill[wash] (axis cs:6.55,-0.5) rectangle (axis cs:7.45,2.05);
\fill[wash] (axis cs:10.55,-0.5) rectangle (axis cs:11.45,2.05);\fill[wash] (axis cs:14.55,-0.5) rectangle (axis cs:15.45,2.05);
\fill[wash] (axis cs:18.55,-0.5) rectangle (axis cs:19.45,2.05);\fill[wash] (axis cs:22.55,-0.5) rectangle (axis cs:23.45,2.05);
\fill[wash] (axis cs:26.55,-0.5) rectangle (axis cs:27.45,2.05);\fill[wash] (axis cs:30.55,-0.5) rectangle (axis cs:31.45,2.05);
\draw[ink!15] (axis cs:-0.6,1.5) -- (axis cs:31.6,1.5);\draw[ink!15] (axis cs:-0.6,0.3) -- (axis cs:31.6,0.3);
\node[anchor=east,align=right] at (axis cs:-1,1.5) {plain vs MTP\\(batch 1, \nmtp)};
\node[anchor=east,align=right] at (axis cs:-1,0.3) {concurrency 1 vs 32\\(plain, \nconc)};
% plain vs MTP: the GDN recurrence (decode kernel against target-verify kernel)
\addplot[only marks,mark=*,mark size=2.6pt,ink,discard if not={kind}{gdn_core},
  nodes near coords={all \cnt: GDN recurrence (decode kernel vs target-verify kernel)},
  every node near coord/.append style={anchor=west,xshift=0.12cm,font=\scriptsize}]
  table[x=layer,y expr=1.5,col sep=comma]{../evidence/state_safety/first_difference_by_module.csv};
% concurrency 1 vs 32: gated RMSNorm of GDN layers, a full-attention layer's output, and the
% MLP down-projection of a batched prefill; the count of prompts above each mark
\addplot[only marks,mark=square*,mark size=2.3pt,ink,discard if not={kind}{gated_norm},
  nodes near coords={\cnt},every node near coord/.append style={font=\scriptsize,anchor=south,yshift=1pt}]
  table[x=layer,y expr=0.3,col sep=comma]{../evidence/state_safety/first_difference_by_module.csv};
\addplot[only marks,mark=triangle*,mark size=2.9pt,ink,discard if not={kind}{attn},
  nodes near coords={\cnt},every node near coord/.append style={font=\scriptsize,anchor=south,yshift=1pt}]
  table[x=layer,y expr=0.3,col sep=comma]{../evidence/state_safety/first_difference_by_module.csv};
\addplot[only marks,mark=diamond*,mark size=2.9pt,ink,discard if not={kind}{mlp_down_proj},
  nodes near coords={\cnt},every node near coord/.append style={font=\scriptsize,anchor=west,xshift=2pt}]
  table[x=layer,y expr=-0.15,col sep=comma]{../evidence/state_safety/first_difference_by_module.csv};
\end{axis}
% key, in absolute centimetres below the axis
\begin{scope}[shift={(0.2cm,-1.45cm)}]
\fill[ink] (-0.07,-0.07) rectangle (0.07,0.07);\node[anchor=west] at (0.15,0) {gated RMSNorm, decode};
\fill[ink] (4.2,0.09) -- (4.12,-0.06) -- (4.28,-0.06) -- cycle;\node[anchor=west] at (4.35,0) {attention output, decode};
\fill[ink] (8.5,0.09) -- (8.59,0) -- (8.5,-0.09) -- (8.41,0) -- cycle;\node[anchor=west] at (8.65,0) {MLP down-projection, batched prefill};
\end{scope}
\end{tikzpicture}}{First differing module output by decoder layer for two configuration pairs.}

\caption{Where equivalent stock configurations first produce different module outputs, by decoder layer. Plain decoding and MTP speculation part at layer 0's GDN recurrence in every prompt, because decode and target verify use different recurrent kernels. Plain decoding at concurrency 1 and 32 (batches of up to 16) part at batch-dependent kernels spread over the depth: the gated RMSNorm of GDN layers (marker size and label give the count), layer 3's attention, or layer 0's MLP in a batched prefill. For plain decoding against MTP, a rerun of 40 prompts that also hashes the caches finds every cache entering the differing recurrence identical; for the concurrency pair these are first differing module outputs only, and whether the caches entering them agree is pending. Each pair's two tapped servers had different pools. Measured with a bit-exact tap~\ev{state-mech,state-cache}.}
\label{fig:divergence}
\end{figure}

\paragraph{How the difference reaches the token} At the position where the runs first choose differently, with competing tokens $a$ and $b$, the analysis compares in each run the order of the two tokens' FP32 accumulator values in the head GEMM, before BF16 rounding. Where a run's BF16 logits differ, rounding is monotone and gives that order. Where they are tied, the float64 dot product of that run's head input with the two rows stands in for the accumulator, but the two can differ by up to $\gamma(A_a+A_b)$, with $A_t$ the sum of absolute products of token $t$'s dot product. A \emph{rounding flip} keeps the accumulator order in both runs, but in one run BF16 rounding merges the two into an exact tie that the index rule resolves against it; an \emph{order flip} reverses the accumulator order between the runs; an \emph{ambiguous} case has a tied run whose float64 gap lies inside the accumulation bound, so its accumulator order cannot be read from the data. The split therefore depends on the error model assumed for the stock accumulator (Table~\ref{tab:divergence}). Where both the BF16 order and a float64 gap outside the bound are available, they never disagree under either model. Under the \modelcons{} most divergences are ambiguous, 93 of 129 and 22 of 40: the stock kernel's possible accumulation error is as large as the gaps that decide these tokens, which is the same fact that makes the gap condition of Theorem~\ref{thm:stock} fail at near ties. Under the \modelhopper{} most plain-versus-MTP divergences are rounding flips (89 of 129), and the concurrency-1-versus-32 divergences divide between order flips (17) and rounding flips (14). Whatever the model, a difference changes the chosen token only where the two leading logits are close: in the concurrency comparison, 151 of its 167 divergences are at an exact BF16 tie, 14 within one BF16 step, and none has a margin above 0.375 nats~\ev{state-noise}.

\begin{table}[htbp]
\small\centering
\begin{tabular}{@{}lrrrlrrr@{}}\toprule
& & \multicolumn{2}{c}{Caused by} & & \multicolumn{3}{c}{Class at the divergence}\\\cmidrule(lr){3-4}\cmidrule(l){6-8}
Pair & Diverged & tie rule & head GEMM & Error model & rounding & order & ambiguous\\\midrule
Plain vs MTP, batch 1 & 129/167 & 0 & 0 & \modelcons & 29 & 7 & 93\\
&&&& \modelhopper & 89 & 16 & 24\\\addlinespace
Plain, concurrency 1 vs 32 & 40/40 & 0 & 0 & \modelcons & 3 & 15 & 22\\
&&&& \modelhopper & 14 & 17 & 9\\\bottomrule
\end{tabular}
\caption{How the first token divergence arises, under the two named models of the stock kernel's FP32 accumulation error: the \modelcons{} ($\gamma=\gcons$), which covers any reduction order, and the \modelhopper{} ($\gamma=\ghop$), Khattak and Mikaitis's blocked \code{wgmma} model with a split-$K$ allowance~\cite{khattak2026}. The tie-rule and head-GEMM columns count divergences caused by the argmax's tie rule on bitwise-equal logits and by the head GEMM on bitwise-equal inputs; they do not depend on the model. The plain-versus-MTP prompts are the 167 that diverged in the concurrency comparison, generated to two tokens past their known divergence; the 40 concurrency prompts are the 16 whose divergence was not an exact tie and 24 of the 151 that were, drawn at random. Measured~\ev{state-mech}.}
\label{tab:divergence}
\end{table}

\paragraph{How often} In the first matrix, at a cap of 16 running requests, plain decoding at concurrency 1 against 32, two passes on one server with identical pools, diverges in 167 of 320 prompts within 256 tokens, 3.42 divergences per 1,000 compared tokens, with a largest margin at the divergence of 0.375 nats. Repeating either configuration on the same server reproduced every token of all 320 prompts, and every log-probability except those of one 128-token prompt, which differ from its prefill onward in both repeats. A fresh server at batch 1 reproduced the first pass bitwise on all 320 prompts, that one included, so its difference on the same server comes from what the prefix cache already held, and the cache-hashing tap places it in the recurrent prefill (below). A fresh server at concurrency 32 reproduced 296 of the 320 outputs; the other 24 diverge at exact ties, 0.35 per 1,000 compared tokens~\ev{state-noise}. The two sessions differed in request arrival, and with it batch composition, and also in their pools: 97,672 against 133,885 attention-cache tokens and 122 against 167 recurrent-state slots~\ev{state-pools}. Plain decoding at batch 1 has an exact BF16 tie between its top two logits at 11.3 of every 1,000 positions and a nonzero gap of at most 0.125 at another 22.2, and no run committed a token that was not its own top-1~\ev{state-noise}. SGLang's deterministic mode, which replaces the default matrix multiplication with a batch-invariant one (DeepGEMM's deterministic BF16 GEMM at the pin's defaults on this GPU, a persistent Triton kernel otherwise)~\cite{sglangdeterministic} and, with the FlashInfer backend, also switches sampling to PyTorch and disables the prefix cache, made plain decoding batch-invariant in an 8-prompt, 128-token check (concurrency 1 against 8: identical tokens and bitwise-identical log-probabilities) but not MTP speculation (2 of 8 prompts diverged). Passing the deterministic KV split to FlashInfer's target-verify plan, as an engine patch (\path{engine/sglang/patches/state/0002-verify-kv-split-deterministic.patch}, behind an environment switch that is off by default and was on for this run), left 28 of 96 MTP prompts diverging between concurrency 1 and 32, 1.55 per 1,000 compared tokens, all at exact ties~\ev{state-noise}. A plausible reason, not tested, is that the draft is not batch-invariant, so accepted lengths, and with them a position's offset inside its verify block, differ between batch sizes.

\paragraph{With pinned pools} A second matrix pins the pools: every server is capped at 8 running requests with 49,152 attention-cache tokens and 40 recurrent-state slots, each was checked to have allocated exactly those, and every pair compares identical pools~\ev{state-pinned,state-pools}. These runs replace the first matrix's comparisons for every pair they cover. A fresh-server repeat reproduces all 320 outputs at concurrency 1 and 316 at concurrency 32, whose 4 divergences are at exact ties (0.06 per 1,000). The drop from 24 comes with both identical pools and the smaller cap, which also narrows the batch compositions, so it is not attributed to the pools alone. Plain decoding at concurrency 1 against 32 on one server diverges in 41 of 320 prompts, 0.63 per 1,000 (39 at exact ties, 2 within one BF16 spacing); with decode batches of at most 8 instead of 16, this rate is not comparable with the 3.42 above. A same-server repeat no longer reproduces every token: 17 of 320 prompts diverge at concurrency 1 (0.25 per 1,000) and 4 at concurrency 32, all at exact ties or within one spacing, where the first matrix's larger pool gave none. No prefill in any pass reused a cached prefix, and one pass computes 126,408 tokens of prompt and output, more than the 49,152-token pool, so the prefix cache evicts during the first pass. The working explanation, not checked at the level of the cache, is the history dependence below: the second pass's requests read different earlier copies of shared prefixes, or run a recurrent prefill that depends on what the prefix cache holds, as one prompt's same-server repeats do. With the prefix cache on, a same-server repeat is therefore not an equality reference. The pool regime alone leaves batch-1 tokens unchanged: the pinned and first-matrix plain runs at concurrency 1 agree on all 320 outputs, and one prompt's log-probabilities differ, from output index 2 and by at most 0.24 nats. With identical pools every MTP configuration, with one, three and five draft steps at top-$k$ 1 and the tree with three steps at top-$k$ 2, diverges from plain decoding at 3.5 to 4.0 per 1,000 compared tokens (164 to 181 of 320 prompts) at both concurrencies, about as often as it diverges from itself between concurrency 1 and 32 (3.1 to 3.5 per 1,000). At concurrency 1 no MTP output is bitwise identical to plain decoding's: the log-probabilities of 313 or 314 of the 320 prompts first differ at output index 1, the first verify forward, and the rest at index 2 or 3, which is consistent with the decode and verify recurrent kernels parting there, although only the tapped prompts locate the kernel. In every pinned pair each divergence is at an exact tie in one run or has both margins within one BF16 spacing, or within 0.25 nats, two spacings, for 11 events in 5 prompts; none falls outside these classes, no run committed a token that was not its own top-1, and log-probability drift without a token change is at most 0.62 nats~\ev{state-pinned}.

\paragraph{Configuration switches with pinned pools} MTP with three steps also ran with the prefix cache off, the overlap scheduler off, deterministic inference or the FP32 head, each at concurrency 1 and 32 on one server with the pinned pools~\ev{state-pinned}. Against the default configuration at concurrency 1, turning the prefix cache off changes the tokens of 98 of the 320 prompts (1.73 divergences per 1,000) and turning the overlap scheduler off those of 161 (3.40); why both change the prefill is below. With deterministic inference MTP still diverges between concurrency 1 and 32, in 168 prompts (3.48 per 1,000, against 3.25 without it), all at exact ties or within one BF16 spacing, so on all 320 prompts deterministic mode does not make MTP batch-invariant, as the 8-prompt check above found. With the FP32 head it diverges there in 130 prompts (2.44 per 1,000); the larger of the two margins is at most 0.19 nats, so the remaining divergences are order flips between near-equal FP32 logits, and the FP32 head removes BF16 ties, not divergences. In the BF16 pairs every divergence is an exact tie, within one spacing, or one of eight near ties on three prompts, where one run's margin is at most 0.375 nats and the other's at most 0.125; no run committed a token that was not its own top-1. The same switches for plain decoding are queued, so a switch's effect on speculation cannot yet be separated from its effect on decoding in general.

\paragraph{Open questions} The plain-decoding switch pairs, the log-probabilities-off control, retraction at concurrency 32, a fresh-server repeat of MTP, the FlashInfer and ReplaySSM decode paths of the recurrent state, and the cache-level check of the concurrency pair are \pend{state-mech}.

\paragraph{History dependence through the prefix cache} Five prompts differ between the tapped and untapped runs, whose servers had different pools~\ev{state-tap,state-pools}. In all five the tapped run first differs from the untapped one in log-probabilities at output index 2, and three of them change tokens, at near ties: at the first changed token the top-two gap is 0 or 0.125 in both runs (\path{tap_check.json}). Untapped runs record no module outputs, so the module-level signature comes from two tapped sessions with different tap versions, which served the same 167 prompts in the same order at batch 1 and also differ for four of the five, in plain decoding and with MTP at three steps (\path{tap_signature.json}). Their pools differ as well, 93,742 against 159,322 attention-cache tokens and 118 against 199 recurrent-state slots for plain decoding, and 84,003 against 166,522 tokens and 55 against 97 slots for MTP; the attention cache never filled in either plain session. In all eight comparisons the first differing module output is layer 3's attention, at output index 2 for plain decoding and at a verify step at output index 3 or 5 for MTP, while its QKV projection and every earlier module output that both tap versions hash agree, which at the same batch shape points to different cached KV. The comparison leaves out the GDN convolution, recurrence and attention outputs, the gated norm and the logits, which the two versions do not hash comparably, so one of those could differ first.

At the pin, after a request's prefill, SGLang inserts its prefix into the prefix cache, matches it again and repoints the request's KV slots to the tree's copies, freeing its own. For tokens the tree already held, such as a chat-template prefix shared with an earlier request, the request then attends over the earlier request's copy of that KV, valid values for the same tokens at the same positions but computed in a prefill of another shape and so not necessarily bitwise equal; under the overlap scheduler decode step 1 is already in flight with the request's own KV, so the switch shows from output index 2 (code reading at the pin). No prefill in the tapped sessions or in the test below reused a cached prefix, and by the code this hybrid cache reuses a prefix only at a stored recurrent state, which these short shared prefixes did not have, so the repoint is the only way one request's KV reaches another.

A deterministic test on stock SGLang serves each of 12 prompts on an empty cache and again right after the earlier prompt that shares the longest prefix with it~\ev{targeted}; the 12 include all five tap-changed prompts. With the prefix cache on, plain decoding and MTP with three steps keep every token but change log-probabilities for 2 of the 12 prompts, \code{mt\_bench-0056} and \code{humaneval-0008}, from output index 2 (and 5 for one MTP prompt); with the prefix cache off all 12 are identical in tokens and log-probabilities, for plain decoding and for MTP with three steps. So with the prefix cache on, a request's output at a fixed configuration and batch shape can depend on which earlier request computed its shared prefix; that is shown for these two prompts. The cache-hashing tap confirms the repoint for both at the level of the cache~\ev{state-cache}. Each prompt was served alone and right after its predecessor, on two fresh servers whose pools differ by at most 30 attention-cache tokens. In both, the first differing module output is layer 3's attention at output index 2; the caches entering decode step 1 are identical, and the keys and values of all 8 attention layers entering step 2 differ at cached positions 3 to 5 for \code{mt\_bench-0056}, whose shared prefix is 6 tokens, and 3 to 21 for \code{humaneval-0008}, with 22 shared tokens, while positions 0 to 2 are bitwise equal in both copies. The other four tap-changed prompts do not change in this test, under plain decoding or MTP. For them attention over a different copy of the cached KV is still the likely mechanism, because of the signature at the same batch shape, but what decides which copy a decode step reads in the two tapped sessions is unknown: the sessions differ in tap version, which changes the host time per forward, in how many tokens the earlier requests generated, and in their pools. One untested candidate is ordering: the repoint is a write issued from the host while the next decode step is already queued, so the step from which the switch shows could depend on timing. The values involved are valid by the code path, no token changed in the history test, and where the tap changed tokens it changed only near-tie decisions; we do not consider it state corruption. The cache-level check of these four prompts is \pend{state-mech}.

\paragraph{History dependence through the recurrent prefill} The cache-hashing tap also found a second route, separate from the repoint~\ev{state-cache}. \code{humaneval-0044}, the 128-token prompt whose log-probabilities differed between same-server repeats in the first matrix, was served five times on one server. With the prefix cache flushed between requests, all 4 repeats are bitwise identical to the first. Without a flush, the second request again prefills all 128 tokens, and its module outputs differ from prompt position 1 on, starting at layer 0's GDN recurrence, with different log-probabilities but the same tokens; the third to fifth reuse a 64-token cached prefix, prefill the other 64, differ from position 64 on and change tokens. That prompt's recurrent prefill therefore depends on what the prefix cache already holds, and the difference shows from the first output token. A related comparison, plain decoding at batch 1 with the prefix cache on and off on 40 prompts, on two servers with different pools (159,324 against 447,027 attention-cache tokens and 199 against 16 recurrent-state slots)~\ev{state-pools}, is identical throughout for the 34 prompts of at most 63 tokens. The 6 prompts longer than one 64-token chunk of the recurrent prefill differ from prompt position 1, at layer 0's GDN recurrence, before any cache is read, and their log-probabilities differ from output index 0; two of them change a token, each at an exact BF16 tie in one run. This is not a question of which copy of the KV a request reads, but it means that a prefix-cache-off control can differ from the prefix-cache-on run from the first output token, as all 6 such prompts tested did. With pinned pools, MTP with three steps at concurrency 1 shows the same on all 320 prompts~\ev{state-pinned}. Turning the prefix cache off changes the top five log-probabilities of the first output token, which come from the prefill alone, in 192 of the 193 prompts longer than 64 tokens and in none of the 127 shorter ones; of the 198 prompts whose outputs differ at all, 192 differ from that token. At the pin, turning the overlap scheduler off also switches the recurrent-state strategy of the prefix cache from \code{extra\_buffer} to \code{no\_buffer}, so it changes two things at once, and it changes the first output token of the same prompts: the overlap-off and prefix-cache-off runs agree at output index 0 on all 320. So the prefill of a prompt longer than one 64-token chunk depends on that strategy, and \code{extra\_buffer} computes it differently from \code{no\_buffer} and from the path without the prefix cache; that the cause is \code{extra\_buffer}'s checkpoint tracking in the prefill kernel is code reading, not verified. The overlap-off and prefix-cache-off runs part later, in 177 prompts at output index 1, the first verify forward, and in 7 further on; the overlap-off server runs SGLang's synchronous speculative path, and which kernel differs there has not been located.

\paragraph{Consequences} A certified head is compared with the stock head at the same batch shape and in the same serving configuration (Section~\ref{sec:same-shape}). And the reference for a speculative arm is stock speculative verification at the same shape, not plain decoding: greedy speculative verification is exact in real arithmetic, but its verify pass runs different kernels from plain decoding, and on every prompt analysed with MTP their outputs already differ at layer 0's recurrence in the first speculative cycle. The DFlash capture is consistent with the same effect at the level of whole outputs: only 39\% of its prompts produced output identical to plain decoding (Appendix~\ref{app:transport-capture}), though its divergences have not been traced to a kernel \pend{drafter-mechanism}. The noise floor applies to fixed-noise sampling in the same way (Section~\ref{sec:decisions}).

\subsection{Recurrent state under speculation}\label{app:state}\label{sec:exactness-state}
Qwen's hybrid architecture makes state correctness part of verification. Reusable state includes attention KV, recurrent tensors, convolution history, positions and backend metadata. With state orientation $S\in\R^{d_v\times d_k}$, the gated-delta update can be written
\begin{equation}
S_t=\alpha_tS_{t-1}+u_tk_t^\top,\qquad u_t=\beta_t\bigl(v_t-\alpha_tS_{t-1}k_t\bigr),\label{eq:gdn}
\end{equation}
which is the Gated DeltaNet rule $S_t=S_{t-1}\alpha_t(I-\beta_tk_tk_t^\top)+\beta_tv_tk_t^\top$~\cite{gdn,deltanet} with the update vector made explicit. The update depends on the preceding state.
\begin{proposition}[Replay of valid derived records]\label{prop:gdn}
If $(\alpha_t,u_t,k_t)$ were recorded along a causal trajectory from $S_0$, applying the first $a$ records to that same $S_0$ reproduces $S_a$ in exact arithmetic:
\[S_a=\Bigl(\prod_{j=1}^{a}\alpha_j\Bigr)S_0+\sum_{t=1}^{a}\Bigl(\prod_{j=t+1}^{a}\alpha_j\Bigr)u_tk_t^\top.\]
\end{proposition}
\begin{proof}
Substitute the expansion at $a-1$ into~\eqref{eq:gdn}; the new gate multiplies every preceding term and the new outer product adds the last summand. Induction completes the proof.
\end{proof}
The qualification about $S_0$ matters: replaying a derived $u_t$ from a different state gives the wrong result even when the token IDs agree. The exact rational reference contains a one-dimensional witness, a record made at state zero and replayed from state two, which gives three where the genuine update gives one~\ev{inherited}. Replaying raw inputs with the true transition is a different operation. Regrouping the recurrence also does not preserve floating-point results (Appendix~\ref{app:contracts}); SGLang's replay fold constrains its operation order and warp count so that it is a bitwise clone of the verify kernel's state update~\cite{gdncode}.

A cycle may emit a correction or bonus token before that token has been consumed into every cache, so the runtime must distinguish the \emph{emitted} prefix from the \emph{materialized} state prefix; treating both as one accepted length invites off-by-one state errors. Publication of speculative state should check that the request's slot and epoch are still current, and memory must not be reused while any stream or captured graph can still write it. An epoch check does not by itself prevent a stale kernel from writing reused memory; ownership and delayed reuse are separate obligations.

\paragraph{What SGLang does, and how it is tested (H5)} During verification SGLang computes the recurrent state and the convolution window after every draft position into per-position buffers and leaves the live state untouched; after acceptance a fused scatter copies the buffers of the last accepted position (for a tree, the last accepted node) into the request's slot (code reading at the pin). The emitted prefix includes the bonus token while the materialized state does not; the bonus token is the first input of the next cycle. Prefix reuse checkpoints the recurrent state every 256 tokens of sequence length into a ping-pong pair of slots, including checkpoints that fall inside a verify cycle, and a freshly allocated slot is cleared by the next forward on the same stream, after any in-flight verification by its previous owner. The MTP draft layer is a full-attention layer, and under greedy verification every committed token is the target's argmax, so an error in draft-side state could lower acceptance but not change the output.

Each point where a speculative cycle touches this state is tested against a reference that isolates it. Truncation by \code{max\_new\_tokens} and stop tokens at every index inside a verify cycle must reproduce the untruncated run's tokens on the same server. Checkpoints taken during speculative decoding, including in the cycle that finished the request, are restored for a continuation and compared with a cold prefill of the same prefix. Requests aborted mid-stream with a recurrent pool of exactly four slots force the next request onto the freed slot while the aborted request's last verification may still be in flight. Chunked prefill at 200 and 256 tokens is compared with an unchunked prefill. And the divergences of every MTP run from plain decoding are bucketed by the previous cycle's commit length, which fixes the draft position that cycle rejected. The first two tests are complete~\ev{targeted}, for MTP with three and five draft steps, top-$k$ 1, the prefix cache, overlap scheduler and CUDA graphs on, one request in flight and the first 40 prompts. Both tests compare output token IDs with the reference prefix; log-probabilities and state are not compared. Truncation leaves every truncated run token-identical to the untruncated run's prefix: 157 of 157 with three steps, keeping 1 to 4 tokens of the final cycle, and 240 of 240 with five steps, keeping 1 to 6. A stop token at every index of a verify cycle also leaves the output token-identical to the untruncated prefix, in 160 of 160 and 240 of 240 cases. In 120 and 200 of them the stop falls inside the draft block, so drafts after it were accepted and folded into the recurrent state before the stop was detected, and none of it changed an output token. Extending each stopped conversation with a new user turn, served warm from the restored checkpoint and cold after a flush, gives token-identical continuations in 139 of 160 and 210 of 240 cases; the rest diverge at exact ties (19 and 30) or within one BF16 step (2, both with three steps), consistent with the warm path's different prefill computation and with the history dependence above.

The same two tests for the tree, with three steps and top-$k$ 2, and for plain decoding also compare the top five log-probabilities~\ev{targeted}. Truncation leaves 159 of 159 tree runs and 39 of 39 plain runs identical in tokens and log-probabilities to the untruncated prefix, and a stop token at every index of a cycle leaves 160 of 160 and 40 of 40 identical up to the stop, 120 of the tree cases with the stop inside the draft block. Warm against cold continuation gives 142 of 160 identical for the tree (17 exact ties, 1 within one spacing) and 33 of 40 for plain decoding (7 exact ties); plain decoding has no speculative state, so such flips occur without it, consistent with the warm path's prefill and the prefix cache's history.

The bucketing by rejection position uses the pinned matrix at concurrency 1 against plain decoding~\ev{state-pinned}. At fragile positions, where plain decoding's top-two gap is at most 0.25 nats, the divergences per fragile position after each commit length of the previous cycle, full acceptance included, lie between 0.061 and 0.098 for MTP with one, three and five steps and the tree, and no length stands out. The rates are compatible with one common rate, with chi-square $p=0.75$, 0.72, 0.54 and 0.38; that is a failure to reject, not a proof of equal rates, and the test treats positions as independent although they cluster by prompt and a prompt's positions stop at its first divergence. No prompt-clustered analysis or bound on a per-length excess rate has been done. No divergence in any of the four MTP configurations falls at a non-fragile position, where a state error would show most plainly. The first verify cycle after the prefill follows no earlier cycle and is counted apart. Its rate is higher for five steps and the tree, 7 of 40 and 6 of 33 fragile positions against 0.077 and 0.081 per fragile position later, and a Fisher exact test against all later cycles gives $p=0.033$ and 0.049 (0.71 and 0.75 for one and three steps). That test is exploratory: it was chosen after looking at the data, it is not corrected for the four configurations and several cycle buckets examined, and the groups differ in more than state, since the first cycle follows the prefill chunk and every prompt contributes it, while later cycles count only prompts that have not yet diverged. It is a lead, not a finding. A test of it is declared before any of its data exist (\path{experiments/state_safety/README.md}, ``Declared follow-up''): on 960 fresh prompts disjoint from the 320, with frozen token IDs, it compares the first cycle with later ones for five steps and the tree against plain decoding at concurrency 1, and uses the same configurations at concurrency 1 against 32, which share the early exit of high-hazard prompts but do not compare verification with plain decoding, as a selection control. A first-cycle excess specific to speculation counts as supported only if the one-sided 95\% bootstrap lower bounds of both the pooled primary log odds ratio and its excess over the control's are above 0; otherwise the result is reported as inconclusive, not as evidence of no effect. Runs that fail the declared provenance, completeness or configuration checks are void, as is the analysis if every prompt is excluded, and a void result reports no statistic~\pend{state-firstcycle}. An effect that also differs between concurrency 1 and 32 would cancel in the control, and a supported result could be a benign difference in numerical path rather than a state error. A check of the recurrent state handed from the prefill to the first verify forward depends on its outcome. For three steps at batch 1 the cache-hashing tap already finds the caches entering the first verify forward, recurrent state included, identical to those entering plain decoding's first decode step, on 40 prompts (Appendix~\ref{app:divergence}); five steps and the tree were not checked. Checkpoint reuse was tested with the prefix cache on~\ev{targeted}. Twelve thinking-mode prompts for MTP with three steps and 8 each for the tree and plain decoding generate 700 tokens, crossing the 256-token checkpoint interval during decode, and prefixes that end at a boundary or 1, 3 or 37 tokens past it are served for 48 tokens warm, after a flush and a regeneration of the original request, and cold, after a flush; tokens are compared, not log-probabilities. A finished request kept only its latest checkpoint, so a warm request restored a decode checkpoint only when its prefix extended past the last boundary the original crossed: in 9 of 100 cases for MTP, 6 of 68 for the tree and 6 of 68 for plain decoding. Where it did, warm and cold gave the same tokens in 9 of 9, 6 of 6 and 4 of 6 cases, the other two at exact ties. A checkpoint taken in the verify cycle that finished the request was restored in every case, and warm and cold agreed in 45 of 53 cases for MTP and 19 of 21 for the tree, the rest at exact ties. Warm and cold reach the boundary state by different computations, decode or verify forwards against a chunked prefill, so bitwise equality is not expected; every difference is an exact tie or within one BF16 step, and plain decoding, which has no speculative state, shows them too, in 9 of 68 cases against 21 of 153 for MTP and 8 of 89 for the tree. No difference points to a wrong restored state, but few decode checkpoints were restored. Aborts with slot reuse, chunked prefill and run-to-run repeats are \pend{state-targeted}.

\subsection{Engine code paths read at the pin}\label{app:integration}
The statements about the stock engine in this paper come from reading SGLang at the pinned commit (Appendix~\ref{app:pins}), a frozen research pin rather than the latest revision. Table~\ref{tab:integration} lists the code paths that decide them and where a certified head would attach; the integration itself is not part of this paper's evidence.

\begin{table}[htbp]
\small\centering
\begin{tabularx}{\linewidth}{@{}P{5.6cm}Y@{}}\toprule
Path under \code{python/sglang/} & Observation\\\midrule
\path{srt/layers/logits_processor.py} & \code{\_compute\_lm\_head} is the single hook through which plain decoding, the MTP drafter and the verifier compute logits: \code{torch.matmul} with BF16 output by default, \code{torch.mm(..., out\_dtype=float32)} under \code{-{}-enable-fp32-lm-head}. A certified head would replace this call for admitted consumers only; log-probability requests and unsupported transforms keep the ordinary path.\\
\path{srt/models/qwen3_5_text.py} & The Qwen3.5 head is tied to \code{embed\_tokens}, which is not a \code{ParallelLMHead}, so SGLang's quantized-head paths do not apply; a low-precision copy must be separate, and the BF16 head stays resident for re-scoring and fallback.\\
\path{srt/models/qwen3_5_mtp.py} & \code{set\_embed\_and\_head} skips the head when embeddings are tied, so a reduced draft head from \code{-{}-speculative-token-map} does not reach the MTP drafter while indices are still remapped (Appendix~\ref{app:levers}).\\
\path{srt/layers/sampler.py} & Unseeded sampling is an exponential race; seeded sampling over the full vocabulary is a Gumbel race keyed by (seed, position, token) in FP64 (\code{multinomial\_with\_seed}), and the certified race uses the same field. The seeded top-$k$ and top-$p$ path keys the same hash by position in the sorted probabilities instead, and the sampler refuses min-$p$ with a seed~\cite{sglangsampler}.\\
\path{srt/models/dflash.py} & DFlash projects draft hidden states through the target head for its candidates; this is where transport's draft-side summaries would be produced.\\
\path{kernels/ops/attention/fla/gdn_replayssm_spec_fold.py} & The replay fold is written as a bitwise clone of the verify kernel and constrains operation order and warp count; its \code{accept\_lens} include the bonus token in that backend's indexing~\cite{gdncode}.\\\bottomrule
\end{tabularx}
\caption{Code paths of the stock engine behind the paper's statements about it, read at the pin; observations, not measurements.}
\label{tab:integration}
\end{table}
